\section{What can presently be inferred}
\label{sec:synthesis}

The calculations establish several affirmative statements. Full vector and STF information can be retained through correlated windows. Exact calibration and common-mode kernels can be exhibited. Finite target missingness has attainable endpoints, and declared selection reweighting has sharp event bounds. Proper rotations and arbitrary orthogonal changes of harmonic coordinates can preserve the same Gram matrix while changing tensor shape. A fixed astrometric Gaussian experiment can support a numerical enclosure of its nested-model statistic and a conservative conditional tail.

They also establish useful limits. A good fit or full design rank does not identify a physical response. The distance-model ranking depends on the assumed group covariance. The supernova covariance improvement does not recover nominal-source parity. Within a restricted history response, the astrometric oracle power is nearly the test size. Split differences cannot constrain arbitrary common foregrounds. A frequency compression cannot recover targets outside its row space. Auxiliary balance cannot produce missing outcomes, and a blank catalogue count cannot be replaced by a guessed value.

These conclusions are compatible. A conditional calculation can answer a precise question while a broader physical question remains unresolved. Its domain and law determine what it answers. In particular, the empirical finite measures used for coupled group fits and astrometric replicates, the constructed CMB reference laws, and the finite Gaussian spin quotient must not be combined into one physical posterior. There is no supplied joint cross-probe likelihood or common physical source state that licenses such a combination.

A physically interpretable continuation needs targeted information. For distance indicators, this includes the accepted-source selection and calibration genealogy, shared peculiar-velocity and calibration covariance, and a controlled optical response on the retained objects. For astrometry, the missing elements are source-class and centroid-to-matter response, true-depth uncertainty, and spatial covariance. For CMB morphology, they are matched processing, split cross-noise, common foreground constraints, and total-field reference laws. An external transfer calculation must be identified by its model, initial data, normalization, and polarization conventions.

Comparing homogeneous anisotropic geometries requires more. A set of scalar scores, an axis, or even an observed full tensor does not determine a Bianchi family. Families can overlap after rotation, projection, source selection, and nuisance response. A family comparison requires physically complete temperature and polarization transfer predictions under a common observation law, together with treatment of family equivalence and model uncertainty \cite{Saadeh2016Framework,Saadeh2016PRL}. No such complete native atlas is supplied by the present calculations.

The report therefore supports a coherent intermediate conclusion: the observed fields and the explicit mathematical obstructions are measurable or derivable at their stated scopes, while their cosmological interpretation remains conditional on missing physical and statistical information. The full tensor records, joint covariance, exact source identities, and distinct probability measures preserve the information needed to pursue that interpretation.
